\documentclass[11pt]{article}

\usepackage[a4paper,margin=1in]{geometry}
\usepackage[T1]{fontenc}
\usepackage{lmodern}
\usepackage{microtype}
\usepackage{amsmath,amssymb,amsthm,mathtools}
\usepackage{booktabs,tabularx}
\usepackage{enumitem}
\usepackage[numbers,sort&compress]{natbib}
\usepackage[colorlinks=true,linkcolor=blue,citecolor=blue,urlcolor=blue]{hyperref}
\usepackage[nameinlink,noabbrev]{cleveref}

\setlength{\emergencystretch}{3em}

\newtheorem{theorem}{Theorem}[section]
\newtheorem{proposition}[theorem]{Proposition}
\newtheorem{corollary}[theorem]{Corollary}
\newtheorem{lemma}[theorem]{Lemma}
\theoremstyle{definition}
\newtheorem{definition}[theorem]{Definition}
\theoremstyle{remark}
\newtheorem{remark}[theorem]{Remark}

\newcommand{\one}{\mathbf 1}
\newcommand{\Lzero}{\mathrm{L0}}
\newcommand{\Lone}{\mathrm{L1}}
\newcommand{\Ltwo}{\mathrm{L2}}

\title{\textbf{Prime-Square Closure of the Frozen
Determinant-Two Logarithmic Shadow}\\
\large Full M\"obius Signs, Fixed Periodic Masks, and the
Growing-Data Firewall}
\author{Liang Wang}
\date{July 2026}

\begin{document}
\maketitle

\begin{abstract}
Let
\[
 D(z)=d+sz,\qquad V(z)=u+az,\qquad su-ad=2
\]
be fixed primitive positive affine forms, and let \(\rho\) be a
fixed bounded periodic sequence.  We prove the unconditional
fixed-data theorem
\[
 \sum_{x/\omega(x)<z\le x}
 \frac{\mu(D(z))\mu(V(z))\rho(z)}{z}
 =o(\log\omega(x))
\]
for every \(1\le\omega(x)\le x\) with \(\omega(x)\to\infty\).
The analytic input is Tao's proved logarithmically averaged
two-point theorem for fixed affine forms.  To close the two
squarefree masks without changing its quantifiers, we use the
Boolean prime-square cutoff
\[
 \mathsf S_P(m)=\prod_{p\le P}(1-\one_{p^2\mid m}).
\]
For fixed \(P\) it has a finite signed CRT expansion, whereas its
full-squarefree error has logarithmic relative mass
\(O(P^{-1})+o(1)\).  The order of limits is essential: \(P\) is
fixed before \(x\to\infty\), and only then tends to infinity.
The result is a genuine fixed-shift arithmetic shadow, but it is
logarithmically weighted and frozen.  It proves no uniformity in
growing coefficients, physical weights or phases, no ordinary
all-prefix power saving, and no positive \(\Ltwo\) input.
\end{abstract}

\section{Fixed determinant-two data}

Fix integers \(a,s\ge1\) and \(d,u\in\mathbb Z\) satisfying
\begin{equation}
 (d,s)=(u,a)=(a,s)=1,\qquad su-ad=2.
\label{eq:primitive}
\end{equation}
After a fixed translation of \(z\), if necessary, assume
\[
 D(z)=d+sz>0,\qquad V(z)=u+az>0
\]
throughout the asymptotic range.  The determinant condition makes
the coefficient pairs \((s,d)\) and \((a,u)\) nonparallel.
TPC-127 identifies these forms with the literal fixed-two
Liouville carrier, while TPC-128 opens their squarefree indicators
by a magnitude cutoff \citep{WangTPC127,WangTPC128}.

Let \(\rho:\mathbb Z\to\mathbb C\) be fixed, periodic of period
\(M\), and satisfy \(|\rho|\le1\).  The words ``fixed'' and
``periodic'' are theorem quantifiers: neither \(M\) nor the values
of \(\rho\) may grow with \(x\) in the theorem below.

\section{A Boolean prime-square cutoff}

Let
\[
 P^\#=\prod_{p\le P}p
\]
and define
\begin{equation}
 \mathsf S_P(m)
 =
 \prod_{p\le P}(1-\one_{p^2\mid m}).
\label{eq:SP-product}
\end{equation}

\begin{lemma}[Exact finite expansion]
\label{lem:expansion}
For every positive integer \(m\),
\begin{equation}
 \mathsf S_P(m)
 =
 \sum_{k\mid P^\#}\mu(k)\one_{k^2\mid m}
 \in\{0,1\}.
\label{eq:SP-expansion}
\end{equation}
Consequently
\begin{align}
 \mathsf S_P(D(z))\mathsf S_P(V(z))
 =
 \sum_{k,\ell\mid P^\#}
 \mu(k)\mu(\ell)
 \one_{k^2\mid D(z)}
 \one_{\ell^2\mid V(z)}.
\label{eq:two-expansion}
\end{align}
\end{lemma}

\begin{proof}
Expand the finite Euler product in \eqref{eq:SP-product}.  Its value
is one exactly when no prime \(p\le P\) has \(p^2\mid m\), and is
zero otherwise.  Multiplying the two one-variable expansions gives
\eqref{eq:two-expansion}.
\end{proof}

For a compatible pair \((k,\ell)\), the two square divisibilities
select one residue modulo \(\operatorname{lcm}(k,\ell)^2\).  In the
determinant-two primitive setting, compatibility in fact forces the
conditions recorded in TPC-128, including \((k,\ell)=1\), but the
following upper bound does not require equality of moduli.

\begin{proposition}[Correct combined modulus]
\label{prop:modulus}
After also resolving the period \(M\), every nonempty component of
\eqref{eq:two-expansion} is a union of residue classes modulo
\[
 L_{M,k,\ell}
 =
 \operatorname{lcm}(M,k^2,\ell^2),
\]
and
\begin{equation}
 L_{M,k,\ell}
 \le
 M\operatorname{lcm}(k,\ell)^2
 \le M(P^\#)^2.
\label{eq:modulus}
\end{equation}
There are at most \(4^{\pi(P)}\) ordered component pairs.
\end{proposition}

\begin{proof}
Intersect the three periodic conditions.  Since
\(\operatorname{lcm}(k^2,\ell^2)=\operatorname{lcm}(k,\ell)^2\),
the first inequality follows from
\(\operatorname{lcm}(M,L)\le ML\).  Both \(k\) and \(\ell\) divide
\(P^\#\), proving the second inequality.  There are
\(2^{\pi(P)}\) divisors of \(P^\#\), independently in each
coordinate.
\end{proof}

\begin{remark}[Why no fourth power appears]
The two divisors \(k,\ell\) are drawn from the same prime set.  The
universal bound is \(M(P^\#)^2\), not \(M(P^\#)^4\).  Writing a
product \(k^2\ell^2\) is harmless only after compatibility has
forced \((k,\ell)=1\); even then \(k\ell\mid P^\#\).
\end{remark}

\begin{lemma}[Residue and content preserve nonparallelity]
\label{lem:reduced-determinant}
Let \(L\ge1\) and \(0\le r<L\).  On the progression
\(z=r+Lt\), put
\[
 c_D=(sL,d+sr),\qquad c_V=(aL,u+ar).
\]
After extracting these positive contents, the two primitive forms
\[
 \widetilde D(t)=\frac{sL}{c_D}t+\frac{d+sr}{c_D},
 \qquad
 \widetilde V(t)=\frac{aL}{c_V}t+\frac{u+ar}{c_V}
\]
have coefficient determinant
\begin{equation}
 \det(\widetilde D,\widetilde V)
 =
 \frac{L(su-ad)}{c_Dc_V}
 =
 \frac{2L}{c_Dc_V}\ne0.
\label{eq:reduced-determinant}
\end{equation}
\end{lemma}

\begin{proof}
The two displayed quotients are primitive by the definitions of
\(c_D,c_V\).  Before content extraction their determinant is
\[
 (sL)(u+ar)-(aL)(d+sr)=L(su-ad)=2L.
\]
Dividing the two coefficient pairs by their contents proves
\eqref{eq:reduced-determinant}.  In particular \(c_Dc_V\mid2L\),
and no CRT residue can make the reduced forms parallel.
\end{proof}

\section{A tail suited to a two-limit theorem}

\begin{lemma}[One-form large-prime-square tail]
\label{lem:tail}
Let \(L(z)=b+cz\) be fixed, primitive and positive on an integer
interval \(I\) of length \(N\), and suppose \(L(z)\le Y\).  Then
\begin{equation}
 \sum_{z\in I}
 |\mu^2(L(z))-\mathsf S_P(L(z))|
 \le
 \sum_{\substack{p>P\\p\le\sqrt Y}}
 \left(\frac{N}{p^2}+1\right)
 \ll_L \frac NP+\sqrt Y.
\label{eq:tail}
\end{equation}
\end{lemma}

\begin{proof}
If \(\mu^2(m)\ne\mathsf S_P(m)\), then \(p^2\mid m\) for some
prime \(p>P\).  For \(p\nmid c\), the condition
\(p^2\mid L(z)\) selects one residue modulo \(p^2\); for
\(p\mid c\), primitivity makes it insoluble.  Summing
\(N/p^2+1\) and using
\(\sum_{n>P}n^{-2}\ll P^{-1}\) proves the claim.
\end{proof}

Unlike the magnitude-truncated divisor sum, \(\mathsf S_P\) is
Boolean.  Therefore
\begin{align}
 &\left|
 \mu^2(D)\mu^2(V)
 -\mathsf S_P(D)\mathsf S_P(V)
 \right|
\notag\\
 &\qquad\le
 |\mu^2(D)-\mathsf S_P(D)|
 +|\mu^2(V)-\mathsf S_P(V)|.
\label{eq:product-tail}
\end{align}
No divisor-function maximum is introduced.

\begin{proposition}[Terminal logarithmic tail]
\label{prop:log-tail}
For the fixed forms above,
\begin{equation}
 \sum_{x/\omega<z\le x}
 \frac{
 |\mu^2(D(z))\mu^2(V(z))
 -\mathsf S_P(D(z))\mathsf S_P(V(z))|}
 {z}
 \ll_{D,V}
 \frac{\log\omega}{P}+1
\label{eq:log-tail}
\end{equation}
uniformly for \(1<\omega\le x\).
\end{proposition}

\begin{proof}
Partition the terminal interval into dyadic pieces
\((T,2T]\), with at most one truncated end piece.  On one such
piece, \cref{lem:tail,eq:product-tail} and \(z^{-1}\le T^{-1}\)
give
\[
 \sum_{T<z\le2T}\frac{\text{tail}(z)}z
 \ll_{D,V}P^{-1}+T^{-1/2}.
\]
There are \(O(1+\log\omega)\) pieces, while their \(T^{-1/2}\)
terms form a geometric sum when ordered from the smallest scale.
The finitely many pieces below a fixed positivity threshold are
absorbed in the constant.
\end{proof}

\section{Full-M\"obius logarithmic closure}

We use the following source theorem only with its native fixed-data
quantifiers.  For fixed positive affine forms
\[
 L_i(z)=a_i z+b_i,\qquad
 a_1b_2-a_2b_1\ne0,
\]
Tao proved
\begin{equation}
 \sum_{x/\omega(x)<z\le x}
 \frac{\lambda(L_1(z))\lambda(L_2(z))}{z}
 =
 o(\log\omega(x))
\label{eq:tao}
\end{equation}
whenever \(1\le\omega(x)\le x\) and
\(\omega(x)\to\infty\) \citep{Tao2016LogChowla}.

\begin{theorem}[Fixed determinant-two full-squarefree closure]
\label{thm:main}
Under \eqref{eq:primitive}, for every fixed bounded periodic
\(\rho\) and every admissible \(\omega\),
\begin{equation}
\boxed{
 \sum_{x/\omega(x)<z\le x}
 \frac{\mu(D(z))\mu(V(z))\rho(z)}{z}
 =
 o(\log\omega(x)).}
\label{eq:main}
\end{equation}
\end{theorem}

\begin{proof}
Fix \(P\).  Insert \eqref{eq:two-expansion} and split \(\rho\) into
its finitely many residue classes.  On each nonempty combined
residue write \(z=r+Lt\).  The two forms become fixed affine forms
in \(t\).  Extract their fixed positive contents.
\Cref{lem:reduced-determinant} gives their exact reduced
determinant \(2L/(c_Dc_V)\), so Tao's nonparallelity hypothesis
survives both the residue split and content extraction.  Complete
multiplicativity extracts fixed Liouville signs.

Discard the finitely many terms with \(t\le0\).  The transformed
terminal interval is
\[
 \frac{x/\omega-r}{L}<t\le\frac{x-r}{L}.
\]
If its lower endpoint is at least one, take the ratio of the two
endpoints; otherwise start at \(t=1\).  In either case this gives
an admissible transformed terminal ratio \(\omega'(x)\to\infty\)
with
\[
 \log\omega'(x)=\log\omega(x)+O_{r,L}(1).
\]
Equation \eqref{eq:tao}, together with
\[
 \frac1{r+Lt}=\frac1L\frac1t+O_{r,L}(t^{-2}),
\]
then gives \(o(\log\omega)\) for every component.  The number of
components is fixed because \(P,M,D,V\) are fixed.

Since \(\mu=\mu^2\lambda\), the difference between the full
M\"obius correlation and this truncated correlation is bounded by
\cref{prop:log-tail}.  After division by \(\log\omega\),
\[
 \limsup_{x\to\infty}
 \frac{|\text{full sum}|}{\log\omega(x)}
 \ll_{D,V}\frac1P.
\]
Now let \(P\to\infty\).
\end{proof}

\begin{remark}[Order of limits]
The theorem does not choose a prescribed \(P=P(x)\).  For each
accuracy request, one first fixes a sufficiently large \(P\), and
then takes \(x\) sufficiently large depending on all fixed data.
This is exactly why \cref{thm:main} is not a growing-family theorem.
\end{remark}

\section{Claim boundary and next gate}

\begin{center}
\small
\begin{tabularx}{\textwidth}{@{}p{0.47\textwidth}X@{}}
\toprule
Statement & Level and status\\
\midrule
Prime-square identity and component bound
& Exact finite theorem (\(\Lzero\)).\\
Large-prime-square and logarithmic tails
& Unconditional elementary theorem (\(\Lzero\)).\\
Full \(\mu(D)\mu(V)\), fixed periodic logarithmic cancellation
& Unconditional frozen arithmetic shadow (\(\Lone\)).\\
Prescribed growing \(P\), affine height, period, or CRT modulus
& Not proved.\\
Physical weight, arbitrary additive phase, ordinary all-prefix
power estimate
& Not proved; actual \(\Ltwo\) gate.\\
H3, \(1/400\), or a prime-pair theorem
& Not proved.\\
\bottomrule
\end{tabularx}
\end{center}

TPC-137 closes a precise gap left by the finite-mask statement of
TPC-129 \citep{WangTPC129}: the full squarefree masks can be returned
for fixed data without claiming a quantitative cutoff rate.  The
next question is whether stronger quantitative results known for
the consecutive shift-one carrier touch the active odd shift-two
carrier.  They do not do so by a parity-preserving relabeling.

\begingroup
\footnotesize
\raggedright
\bibliographystyle{plainnat}
\bibliography{references}
\endgroup

\end{document}
